\documentclass[12pt]{article}
\usepackage[a4paper,margin=1in]{geometry}
\usepackage[T1]{fontenc}
\usepackage{lmodern,microtype}
\usepackage{amsmath,amssymb,amsthm,mathtools}
\usepackage{booktabs,array,enumitem,xcolor}
\usepackage[numbers,sort&compress]{natbib}
\usepackage[colorlinks=true,linkcolor=blue!55!black,citecolor=blue!55!black,urlcolor=blue!55!black]{hyperref}
\linespread{1.04}

\newtheorem{theorem}{Theorem}[section]
\newtheorem{proposition}[theorem]{Proposition}
\newtheorem{corollary}[theorem]{Corollary}
\theoremstyle{remark}
\newtheorem{remark}[theorem]{Remark}

\title{Clock Rank, Cycle Length, and Cloud Degree\\
\large An Exact Integer Translation and a Seven-Level Calibration Audit}
\author{Liang Wang}
\date{July 2026}

\begin{document}
\maketitle

\begin{abstract}
The finite-cycle route introduces three integer scales: the RH-82 reset clock
rank $r_\sigma$, the RH-15 effective cloud degree $N_\sigma$, and the cycle
length $L_\sigma=N_\sigma+1$.  We derive their exact bookkeeping relation.
Let
\[
 h_\sigma=\frac{\log(1/\sigma)}{2\log\lambda},
 \qquad r_\sigma=\lceil h_\sigma\rceil+2,
 \qquad d_\sigma=h_\sigma-N_\sigma.
\]
Then
\[
 r_\sigma-L_\sigma=\lceil d_\sigma\rceil+1.
\]
Thus every bounded cloud-degree defect immediately becomes a bounded finite
rank-to-cycle correction.

Joining the seven RH-15 cloud rows to the half-log clock gives
$1.4456\le d_\sigma\le2.3376$ and therefore
\[
 r_\sigma-L_\sigma\in\{3,4\}.
\]
Four rows have gap three and three rows gap four.  At the only scale shared
with the certified RH-151 reset atlas, $\sigma=0.01$, the actual reset rank is
seven, the cloud degree is three, the cycle length is four, and the gap is
three, exactly as the translation predicts.

The six smaller-noise ranks are formal RH-82 clock values, not observed
RH-151 reset packets.  One physical overlap cannot select a unique correction
$L=r-3$ versus $L=r-4$, and seven cloud fits cannot prove an asymptotic degree
law.  The result supplies a precise calibration corridor, not a physical
cycle theorem.  A valid next step must derive the cycle from time-domain or
transfer data and then test which integer branch survives.
\end{abstract}

\section{Three clocks that must not be identified by notation}

The prime-dynamics program now contains three different finite dimensions.
First, RH-82 proposes the half-log source packet clock
\citep{WangRH82}.  Second, RH-15 selects an effective outer cloud degree from
the noisy transfer spectrum \citep{WangRH15}.  Third, RH-176--177 realize a
degree-$N$ geometric factor by a reduced cycle of length $L=N+1$
\citep{WangRH176,WangRH177}.

These dimensions are related but not equal by definition:
\begin{description}[leftmargin=3.5em,style=nextline]
\item[$r_\sigma$] source-memory packet rank;
\item[$N_\sigma$] fitted degree per geometric cloud copy;
\item[$L_\sigma$] full cycle length, including the removed stationary mode.
\end{description}
Any physical bridge must explain their offsets rather than suppress them.

\section{Half-log clock and exact translation}

Fix the deterministic expansion parameter
\[
 \lambda=1.6785735104283224
\]
and define
\begin{equation}\label{eq:h}
 h_\sigma=\frac{\log(1/\sigma)}{2\log\lambda}.
\end{equation}
The RH-82 clock rank with offset two is
\begin{equation}\label{eq:r}
 r_\sigma=\lceil h_\sigma\rceil+2.
\end{equation}
For an observed cloud degree $N_\sigma$, define
\begin{equation}\label{eq:Ld}
 L_\sigma=N_\sigma+1,
 \qquad
 d_\sigma=h_\sigma-N_\sigma.
\end{equation}

\begin{theorem}[Clock--cycle integer translation]\label{thm:translation}
For every real $h>0$ and integer $N\ge0$, if
$r=\lceil h\rceil+2$, $L=N+1$, and $d=h-N$, then
\begin{equation}\label{eq:translation}
 r-L=\lceil d\rceil+1.
\end{equation}
More generally, a rank offset $a\ge0$ gives
$r-L=\lceil d\rceil+a-1$.
\end{theorem}

\begin{proof}
Since $N$ is an integer,
$\lceil h\rceil-N=\lceil h-N\rceil=\lceil d\rceil$.  Therefore
\[
 r-L=(\lceil h\rceil+2)-(N+1)
 =\lceil d\rceil+1.
\]
The offset-$a$ formula is identical.
\end{proof}

\begin{corollary}[Defect corridor]
If $1<d\le2$, then $r-L=3$.  If $2<d\le3$, then $r-L=4$.  Hence a cloud
degree staying between one and three units below the half-log horizon yields
a cycle length three or four units below the offset-two packet rank.
\end{corollary}

The result is integer arithmetic, not a dynamical theorem.  Its value is that
every empirical degree discrepancy can now be translated without ambiguity.

\section{Seven-level joined audit}

The RH-15 cloud table contains seven noise scales.  For each row we recompute
$h_\sigma$ from \eqref{eq:h}, read $N_\sigma$, set $L=N+1$, and evaluate the
formal clock rank \eqref{eq:r}.

\begin{center}
\small
\begin{tabular}{rrrrrrr}
\toprule
$\sigma$ & $h_\sigma$ & $N$ & $L$ & formal $r$ & $d=h-N$ & $r-L$\\
\midrule
$10^{-2}$ & 4.4456 & 3 & 4 & 7 & 1.4456 & 3\\
$4\cdot10^{-3}$ & 5.3302 & 3 & 4 & 8 & 2.3302 & 4\\
$2\cdot10^{-3}$ & 5.9993 & 4 & 5 & 8 & 1.9993 & 3\\
$10^{-3}$ & 6.6684 & 5 & 6 & 9 & 1.6684 & 3\\
$5\cdot10^{-4}$ & 7.3376 & 5 & 6 & 10 & 2.3376 & 4\\
$2\cdot10^{-4}$ & 8.2221 & 6 & 7 & 11 & 2.2221 & 4\\
$10^{-4}$ & 8.8912 & 7 & 8 & 11 & 1.8912 & 3\\
\bottomrule
\end{tabular}
\end{center}

The archived column called ``degree defect from half horizon'' agrees with
$d=h-N$ in every row.  The exact translation identity has zero failures.
The gap sequence is
\[
 3,4,3,3,4,4,3.
\]
This bounded offset is more informative than a raw correlation: it states the
precise finite integer ambiguity that a physical construction must resolve.

\section{The sole physical reset overlap}

The RH-151 reset atlas covers
$\sigma=0.16,0.08,0.04,0.02,0.01$ \citep{WangRH151}, while RH-15 covers
$0.01$ and smaller noise.  Their intersection is the single scale
$\sigma=0.01$.  At that scale:
\[
 h=4.4456227,
 \qquad
 r_{\rm reset}=7,
 \qquad
 N_{\rm cloud}=3,
 \qquad
 L_{\rm cycle}=4.
\]
Thus
\[
 r_{\rm reset}-L_{\rm cycle}=3,
\]
and the actual reset rank agrees with the formal RH-82 clock.

This one overlap is a consistency check, not a fitted law.  In particular,
it cannot distinguish between a rule that usually uses gap three and one that
switches between gaps three and four according to a secondary phase or parity
condition.

\section{Rounding sensitivity of the ceiling clock}

The integer translation is exact once $h_\sigma$ is fixed, but the map
$h\mapsto\lceil h\rceil$ is discontinuous at integers.  Define the ceiling
margin
\[
 m_\sigma=\operatorname{dist}(h_\sigma,\mathbb Z).
\]
If an uncertainty interval for $h_\sigma$ has radius smaller than
$m_\sigma$, the formal rank is stable.  If it crosses an integer, the clock
rank and rank--cycle gap can change by one.

Among the seven rows, $\sigma=0.002$ is exceptionally close to a boundary:
\[
 h_{0.002}=5.9993011961,
 \qquad
 m_{0.002}=6.9880\times10^{-4}.
\]
A perturbation of the half-clock larger than this can move the ceiling from
six to seven and change the formal packet rank from eight to nine.  The next
smallest margin in the table is about $0.1088$ at $\sigma=10^{-4}$; the other
rows are much less sensitive.

\begin{proposition}[Interval-stable translation]
Let $h\in[h_-,h_+]$ and fix an integer degree $N$.  If
$\lceil h_-\rceil=\lceil h_+\rceil$, then the clock rank and rank--cycle gap
are constant throughout the interval.  Otherwise the possible gaps are the
consecutive integers
\[
 \lceil h_-\rceil-N+1,
 \ldots,
 \lceil h_+\rceil-N+1.
\]
\end{proposition}

\begin{proof}
The cycle length $N+1$ is fixed, and the only varying integer is
$\lceil h\rceil$.  The ceiling takes every consecutive integer between its
endpoint values.
\end{proof}

For a validated all-level clock, uncertainty in the expansion parameter
$\lambda$ must therefore be propagated through $h_\sigma$ before rounding.
The current audit uses the fixed archived value exactly as specified and does
not attach an interval to it.

\section{What could generate the offset?}

The integer difference has several possible structural sources:
\begin{enumerate}[leftmargin=2em]
\item the reset packet includes the explicit rank offset two used to protect
weak modes, while the reduced cycle removes one stationary mode;
\item the cloud degree is selected by a radial/phase threshold and can lag the
half-log horizon by one additional mode;
\item left and right packet channels may share or split stationary and parity
directions differently from the two cloud copies;
\item finite-resolution eigensolvers may undercount a mode near the separating
radius.
\end{enumerate}
The arithmetic theorem does not choose among these mechanisms.  A physical
cycle construction must expose which dimensions are stationary, protected,
or genuinely resonant.

\section{Two calibration branches}

The observed corridor suggests two finite hypotheses:
\begin{align}
 \mathsf C_3:&\quad L_\sigma=r_\sigma-3,
 \label{eq:c3}\\
 \mathsf C_4:&\quad L_\sigma=r_\sigma-4.
 \label{eq:c4}
\end{align}
Neither is asserted globally.  The current seven rows select $\mathsf C_3$
four times and $\mathsf C_4$ three times.  A third possibility is a switching
law determined by the fractional part of $h_\sigma$ or by an independently
measured boundary phase.

\section{Total mode count and two falsifiable branches}

The doubled reduced cycle has total nonstationary rank
\begin{equation}\label{eq:mode-count}
 m_\sigma=2N_\sigma=2(L_\sigma-1).
\end{equation}
Combining this with a fixed rank--cycle gap $g=r-L$ gives an exact conversion
from reset rank to cloud count.

\begin{proposition}[Mode counts under the two calibration branches]
\label{prop:mode-branches}
If $L=r-3$, then
\begin{equation}\label{eq:m3}
 m=2r-8.
\end{equation}
If $L=r-4$, then
\begin{equation}\label{eq:m4}
 m=2r-10.
\end{equation}
At the physical overlap $r=7$, the observed cloud count six selects the first
formula and rejects the second at that scale.
\end{proposition}

\begin{proof}
Substitute $N=L-1=r-g-1$ into $m=2N$.  For $g=3$ and $g=4$ this gives
\eqref{eq:m3} and \eqref{eq:m4}.  When $r=7$, the values are six and four,
respectively; RH-15 records six cloud eigenvalues.
\end{proof}

This is stronger than comparing only $L$ and $r$: it predicts the full doubled
cloud multiplicity.  At a future scale where both a certified reset rank and
a validated Riesz cloud count are available, the two branches can be tested
without fitting phases or radii.  If neither count matches, the simple fixed
offset model is false even if the half-log and cloud degree remain correlated.

The current seven formal rows give predicted total counts
\[
 6,6,8,10,10,12,14,
\]
exactly the archived values $2N$.  This agreement is tautological on the
cloud side because $m=2N$ defines the model count; only the single reset
overlap tests a cross-family prediction.

\section{Falsification criteria for the calibration corridor}

The finite corridor should be abandoned or revised if any of the following
occurs:
\begin{enumerate}[leftmargin=2em]
\item a validated cloud degree has $h-N\notin(1,3]$ persistently, producing a
gap outside $\{3,4\}$;
\item two or more scales with certified reset packets disagree with the
formal clock rank after interval propagation;
\item the physical Riesz cloud has a mode count different from $2N$ because
one copy is absent, extra stationary modes survive, or multiplicities split;
\item a time-domain cyclic construction selects a length incompatible with
both $r-3$ and $r-4$.
\end{enumerate}
These are clean negative outcomes.  The purpose of the integer dictionary is
to make the cycle hypothesis testable before a large determinant calculation,
not to protect it from contradictory data.

\begin{proposition}[What would make calibration rigorous]
Suppose an all-level cloud theorem proves integers $N_\sigma$ satisfying
\[
 a<h_\sigma-N_\sigma\le b
\]
for fixed real $a,b$.  Then the rank--cycle gap lies in the finite set
\[
 \{\lceil a\rceil+1,\lceil a\rceil+2,
 \ldots,\lceil b\rceil+1\},
\]
with endpoint adjustments when $a$ is an integer.  If the defect eventually
lies in one unit interval $(k-1,k]$, then the gap is eventually the single
integer $k+1$.
\end{proposition}

\begin{proof}
Apply Theorem~\ref{thm:translation} and monotonicity of the ceiling function.
\end{proof}

Thus the needed analytic target is not exact equality of two ranks.  It is an
eventual unit-width bound on the cloud-degree defect.

\section{Audit boundary}

The joined script reads the archived RH-15 CSV and RH-151 JSON, recomputes the
RH-82 half-log clock, and checks the integer identity.  It has seven rows, one
actual reset overlap, zero translation failures, and zero reset-clock
mismatches on that overlap.

The six smaller-noise values $r=8,8,9,10,11,11$ are predictions of the clock
formula only.  No reset packets were computed or certified at those six
scales in this paper.  Likewise, the RH-15 values $N=3,4,5,6,7$ are numerical
cloud selections, not an all-level theorem.

\section{Route consequence}

The finite-cycle branch now has an exact integer dictionary and a narrow
empirical corridor:
\[
 \text{cloud degree }N
 \leftrightarrow
 \text{cycle length }L=N+1,
 \qquad
 r-L\in\{3,4\}\text{ on seven rows}.
\]
The next step cannot be more bookkeeping.  It must place the cycle operator
inside, or close to, a physical transfer-space block and use explicit root
spacing to construct Riesz shells.  RH-180 gives that conditional theorem.

No unique cycle calibration, asymptotic degree law, physical Riesz cloud,
Gate-A determinant, Hilbert--Polya operator, zeta identity, or Riemann
Hypothesis is proved.

\bibliographystyle{plainnat}
\bibliography{references}
\end{document}
